% !TeX root = ../main.tex
% Add the above to each chapter to make compiling the PDF easier in some editors.

\chapter{Methodology}\label{chapter:methodology}


\section{Problem Statement and Prediction Target}

In this thesis, we address the transfer of the T3 runtime prediction approach from Umbra to PostgreSQL.
The central methodological goal is to preserve T3’s good performance~\parencite{t3}.
Thus, we use its main concepts and adapt plan conversion and feature extraction to PostgreSQL-specific plan information accordingly.

The final prediction target is total query runtime.
Each query plan is first decomposed into pipelines and a feature vector is computed for each pipeline.
These per-pipeline feature vectors are then summed into a single query-level feature vector,
from which the model predicts the total query runtime.
The pipeline decomposition is therefore used for feature extraction,
not as a separate prediction unit.

\section{Data and Preprocessing}
Regarding the training data, we use the public ZeroShot benchmark data introduced with the ZeroShot Cost Model.
It is a PostgreSQL benchmark, which makes it directly suitable for our implementation.
Furthermore, it consists of real-world datasets (e.g., IMDB) and synthetically generated workloads designed to cover diverse query structures.
This includes \emph{SPAJ} (Select-Project-Aggregate-Join) queries with conjunctive predicates.
Overall, the benchmark contains 20 databases with large variation in tables, columns, and foreign-key relationships \parencite{zeroshot}.

Having this multitude of databases allows us to test the generalization capabilities of our model.
Similar to the T3 paper~\parencite{t3}, we can train the model on all but one database and test on the remaining one.
Moreover, as the widely used JOB workload is included, we can compare our performance predictions to other implementations later on.

Following the benchmark construction described in the ZeroShot paper, queries are generated per database and executed on a PostgreSQL DBMS.
Then, they are stored as raw traces (i. e., the output of \texttt{EXPLAIN ANALYZE}).
Next, the authors preprocess the queries to remove very long-running queries (queries above 30\,s are excluded) to keep the benchmark practical and reproducible.
A further preprocessing step is to parse the raw traces into structured plan representations.
These parsed plans contain the operator nodes found in the raw traces, but organize them in a tree structure and gather node-wise attributes.
For every node, e.g., operator type, estimated/actual cardinalities, widths, and runtimes are stored~\parencite{zeroshot}.

The tree structure can be used to create pipelines and the node information will be used to train the model.
Consequently, the parsed plans offer a good starting point for our pipeline and we can begin without further preprocessing.



\section{Plan Representation}

Each parsed plan is a rooted tree of operator nodes, where every node carries per-operator attributes such as cardinalities, costs, widths, and timings.

We convert this tree recursively into the same nested operator structure that the original T3 implementation expects for Umbra plans~\parencite{t3}.
PostgreSQL operator names are mapped to a small set of abstract operator types that T3 uses internally for pipeline assignment.

For joins, the child order is normalized so that the build subtree becomes the left child and the probe subtree the right child, matching T3's pipeline rules.

Since the feature vector uses PostgreSQL-native attributes, the original per-node attributes are preserved alongside the converted abstract structure.
The abstract operator types serve only for pipeline assignment; feature extraction reads directly from the original PostgreSQL-native attributes.

\section{Pipeline Decomposition}

Umbra plans group operators into \emph{pipelines}: sets of operators that can produce tuples in a pipelined fashion, separated by \emph{pipeline breakers}. 
Namely, these are operators/stages where data is materialized. For instance, sorting, grouping, or materialization, but also the build side of a hash join~\parencite{t3}.

Since PostgreSQL plans do not provide such a pipeline breakdown, we need to derive it from the converted tree.
We assign each node to a pipeline by walking the converted tree with the same high-level rules as in the original T3 implementation~\parencite{t3}.
Thus, PG's sort and aggregate nodes break pipelines; a PG-materialize node is treated as a breaker.
Joins place the build side in one pipeline and keep the join together with the probe side in the other pipeline.

\section{Feature Vector Design}\label{sec:methodology-implementation}

We retain the pipeline decomposition for feature extraction rather than opting for a complete query-level feature vector for three reasons.

By mirroring T3's pipeline-based feature design~\parencite{t3},
we can use the same pipeline-centric features such as operator percentages, which are defined within a pipeline boundary.
Further, following T3's ablation study, which includes per-pipeline predictions, we can compare prediction granularities more thoroughly in Section~\ref{sec:evaluation_ablation}.
Finally, PostgreSQL-specific issues affect the per-pipeline and per-tuple approaches
in our setting, as discussed in the evaluation (Section~\ref{sec:evaluation_ablation}).
The decomposition keeps the model extensible so that future work can address those issues.

The original T3 feature mapper for Umbra builds a feature vector from execution phases within each pipeline (e.g.\ scan, build, probe phases), as well as from cardinalities and tuple sizes attached to those phases~\parencite{t3}.
Since ZeroShot plans contain PostgreSQL-specific data, we construct the feature vector for each pipeline differently, after the pipeline decomposition described above.

Once the feature vector is constructed, we can train the model.
Our feature vector differs from the original T3 feature vector in the following ways:



\begin{enumerate}
  \item \textbf{PostgreSQL-Native Operator Features.} The ZeroShot data exposes PostgreSQL-native operator nodes. 
  Thus, we add entries for all occurring operators in the feature vector.
  
  
  \item \textbf{PostgreSQL Planner Cost Estimates.} We include sums and extrema of PostgreSQL's cost estimates in the feature vector.
  They are optimizer-internal estimates of work and startup cost per node and are an additional PostgreSQL-specific signal that is not available in the Umbra plans.
  \item \textbf{Planned Parallelism.} Another PostgreSQL-specific signal is the planned parallelism. Thus, we take the sum of planned workers over all operators in the pipeline.
  This provides a dedicated parallelism signal in our feature vector.
  \item \textbf{Sizes.} T3 uses the size of output tuples in bytes for materializing operators~\parencite{t3}.
  In Umbra, T3 derives the tuple size by summing the per-attribute byte sizes from the operator's IU (information unit) list,
  giving a precise estimate of the tuple byte width.
  PostgreSQL provides a direct semantic equivalent: the estimated average output row width in bytes per node (est\_width).
  We therefore use it as the tuple size for each operator.
  The main difference to the Umbra approach is that it is a pre-aggregated planner estimate~\parencite{postgresql_docs_explain}
  rather than a precise sum of per-attribute type metadata.
\end{enumerate}

Regarding stages, T3 derives stage assignments (Scan, Build, Probe, PassThrough) from Umbra-specific operator types and their position within the pipeline.
In PostgreSQL, however, this information is already encoded in the node type itself.
Umbra represents both phases of a hash join as a single HashJoin operator appearing in two separate pipelines with different stages.
In contrast, PostgreSQL exposes them as two distinct plan nodes -- a Hash node for the build side and a Hash Join node for the probe side.
A stage feature is therefore redundant in the PostgreSQL representation, and we omit it accordingly.
Thus, we use operator-level features.

Furthermore, we add pipeline-level features.
They produce additional metadata in a compact manner, e.g., by using averages and sums over operators.
For instance, we consider planned parallelism, width, and cost estimates per pipeline and not per operator.
Thus, the feature vector stays more compact.

Regarding cardinalities, either actual or estimated cardinalities are used throughout the feature vector, depending on the training and inference configuration.
Full details of the feature vector layout are provided in Chapter~\ref{chapter:implementation}.


\section{Training}

We keep T3's LightGBM decision tree setup~\parencite{t3} and use the newly constructed vectors.
A main difference is that we predict the total query runtime directly,
rather than predicting individual pipeline durations.

The training label is the total observed query runtime.
The model input is formed by summing all per-pipeline feature vectors of a query
into a single query-level feature vector.
This per-query prediction approach was found to perform best for our PostgreSQL setting,
as demonstrated in the ablation study (see Section~\ref{sec:evaluation_ablation}).

In T3~\parencite{t3}, the training target is the per-tuple execution time for each individual pipeline instead,
and the predicted pipeline times (gathered from the per-tuple predictions) are summed to obtain the query runtime.
We evaluate all three granularities -- per-query, per-pipeline, and per-tuple --
in the ablation study and discuss why per-query works best for our case.

Our model uses gradient boosted decision trees as implemented in LightGBM~\parencite{lightgbm}.
The booster is trained with a mean absolute percentage error (MAPE) objective on the training rows,
in line with the original T3 work~\parencite{t3}.

Furthermore, we use the same training parameters as in the original T3 work~\parencite{t3}.
Hence, the prediction targets $t$ are transformed as $-\log(t)$ to stabilize training.
Also following T3, we cap the number of boosting rounds at 200 trees to mitigate overfitting.

Regarding the training data, we mirror T3's train/validation split.
We sample 80\% of the training data for training and 20\% for validation~\parencite{t3}.
The default seed used for shuffling the training data is 42.

Concerning the train/test split, we match the training intention of the ZeroShot benchmark~\parencite{zeroshot} and of T3's cross-instance evaluation~\parencite{t3}.
Both aim to measure whether the trained model transfers to an unseen database instance.
Thus, we use a leave-one-out split on the ZeroShot data:
one instance is held out as the test set and all other instances form the training data.


